%=============================================================================!
% 11.2  VELOCITIES TOO: RATTLE                                                !
%=============================================================================!
\section{Velocities too: RATTLE}
\label{sec:cs-rattle}

SHAKE corrects positions, and the Störmer--Verlet step it was built on
carries no velocities. The loop of \cref{sec:md-loop} uses velocity
Verlet, which needs the velocities at the same times as the positions,
for the kinetic energy and to start the next step. Andersen's RATTLE
gives velocity Verlet the same constraints\cite{Andersen1983}; this
section derives it, checks it against an independent code, and counts
the motions that are left.

\subsection*{Two corrections}

Velocity Verlet is a half kick, a drift and a half kick
(\cref{sec:in-splitting}). With the constraint force added to the first
half kick, the step of atom $i$ begins
\begin{align*}
   \vb{v}_i(t + \tfrac12\dt) &= \vb{v}_i(t) + \frac{\dt}{2m_i}
      \bigl(\vb{F}_i + \vb{F}_i^{\mathrm c}\bigr) ,\\
   \vb{r}_i(t + \dt) &= \vb{r}_i(t) + \dt\,\vb{v}_i(t + \tfrac12\dt)
   = \tilde{\vb{r}}_i + \frac{\dt^2}{2m_i}\,\vb{F}_i^{\mathrm c} ,
\end{align*}
with the provisional position now $\tilde{\vb{r}}_i = \vb{r}_i(t) +
\dt\,\vb{v}_i(t) + (\dt^2/2m_i)\vb{F}_i$. This has the form of the
constrained Störmer step with $\dt^2/2$ in place of $\dt^2$, a factor
that $\eta = \dt^2\lambda/(2m_{\mathrm r})$ absorbs, so \texttt{shake}
corrects the positions unchanged. The half-step velocity drove the drift,
so it must take the same correction, divided by $\dt$; equivalently,
$\vb{v}_i(t + \frac12\dt) = [\vb{r}_i(t + \dt) - \vb{r}_i(t)]/\dt$ once the
positions are corrected.

The second half kick, with the forces at the new positions, leaves
velocities that would in general stretch or shrink the bonds. The length
of a bond stays $d$ only if it is not changing, and, differentiating
component by component as for $\lvert\vb{v}\rvert^2$ in
\cref{sec:mo-circle},
\begin{equation*}
   \frac{\dd}{\dd t}\,\tfrac12\lvert\vb{r}_{ij}\rvert^2
   = \vb{r}_{ij}\cdot\vb{v}_{ij} ,
   \qquad \vb{v}_{ij} = \vb{v}_j - \vb{v}_i ,
\end{equation*}
which must be zero at $t + \dt$: the two atoms may move side by side or
turn about each other, but not approach or separate along the bond. A
second constraint force, along the new bond $\vb{r}_{ij}$, changes $\vb{v}_j$
by $(m_{\mathrm r}\zeta/m_j)\,\vb{r}_{ij}$ and $\vb{v}_i$ by
$-(m_{\mathrm r}\zeta/m_i)\,\vb{r}_{ij}$, for some number $\zeta$, so that,
as for the positions, the relative velocity changes by $m_{\mathrm
r}\zeta(1/m_j + 1/m_i)\vb{r}_{ij} = \zeta\vb{r}_{ij}$. Then
\begin{align*}
   (\vb{v}_{ij} + \zeta\vb{r}_{ij})\cdot\vb{r}_{ij} &= 0 ,\\
   \zeta &= -\frac{\vb{v}_{ij}\cdot\vb{r}_{ij}}{d^2}
   && \text{as $\lvert\vb{r}_{ij}\rvert^2 = d^2$.}
\end{align*}
This condition is linear in $\zeta$, so one correction makes a single
bond exact, and bonds that share atoms are swept until all hold, as in
SHAKE. For atom $i$ of mass 1 and atom $j$ of mass 3, joined along $x$
with $d = 1$, so that $\vb{r}_{ij} = (1, 0, 0)$, and with $\vb{v}_i =
\vb{0}$ and $\vb{v}_j = (0.3, 0.4, 0)$, $\zeta = -0.3$, and the corrected
velocities are $(0.225, 0, 0)$ and $(0.225, 0.4, 0)$: the relative motion
along the bond is gone, the sideways motion of atom $j$ is untouched.

Every correction, of positions or of velocities, changes $m_j\vb{v}_j$ by
$m_{\mathrm r}\zeta\vb{r}_{ij}$ and $m_i\vb{v}_i$ by the opposite (or the
positions likewise), so the total momentum and the centre of mass are
unchanged by the constraints, as \cref{sec:ne-bodies} requires of forces
between the atoms of a system. In the example, the momentum along $x$ is
$3\times0.3 = 0.9$ before and $0.225 + 3\times0.225 = 0.9$ after, and
along $y$ it stays $3\times0.4 = 1.2$.

The function \texttt{rattle\_step} of \texttt{mdlab.constraints} is the
whole step:
\begin{code}
    v_half = v + 0.5 * dt * force_to_accel * np.asarray(forces) / m
    r, _ = shake(r0 + dt * v_half, r0, m[:, 0], bonds, lengths, **keep)
    v_half = (r - r0) / dt
    u, f = model(r)
    v = v_half + 0.5 * dt * force_to_accel * f / m
    v, _ = rattle_velocities(r, v, m[:, 0], bonds, lengths, **keep)
    return r, v, u, f
\end{code}
\noindent with \texttt{force\_to\_accel} the factor of \cref{sec:ne-atoms}
that turns a force in \si{\electronvolt\per\angstrom} over a mass in amu
into an acceleration in \si{\angstrom\per\femto\second\squared}, and
\texttt{keep} passing on the cell and the tolerance. The model is called
once per step, as in velocity Verlet; the constraints add only the sweeps,
and for the water box of \cref{sec:cs-water} one SHAKE to a tolerance of
$10^{-10}$ after a step of \SI{2}{\femto\second} takes about a tenth of the
time of one call of the model.

Running the step forwards and then backwards from the reversed velocities
retraces it. Fifty steps of \SI{2}{\femto\second} of the water box of
\cref{sec:cs-water}, out and back with a tolerance of $10^{-14}$, recover
the starting positions to $\num{5.2e-12}$\,\AA\ and the velocities to
$\num{1.9e-13}$\,\AA/fs. Leimkuhler and Skeel showed that RATTLE, like
velocity Verlet, is symplectic for motion confined to the
constraints\cite{LeimkuhlerSkeel1994}, so its energy should stay in a
band, as \cref{sec:cs-water} measures.

\subsection*{Checked against ASE}

ASE carries out the same two corrections: its \texttt{VelocityVerlet},
given the constraint \texttt{FixBondLengths}, corrects the positions and
half-step velocities after the drift and the velocities after the second
half kick, with the update of \cref{eq:cs-linear} for the positions and
the $\zeta$ above for the velocities. Given the forces of the
\texttt{mdlab} water model, the same start and the same femtosecond
(\cref{sec:md-loop}), 50 steps of \SI{1}{\femto\second} of the box of 64
rigid molecules agree to $\num{2.1e-12}$\,\AA\ in the positions and
$\num{1.5e-13}$\,\AA/fs in the velocities, with both codes' bonds held to
$\num{1.5e-13}$\,\AA\ at a tolerance of $10^{-13}$.

\subsection*{The motions that are left}

Each independent distance constraint removes one way of moving. A system
of $N$ atoms with $n_{\mathrm c}$ independent held distances has $3N - n_{\mathrm c}$ degrees of
freedom (\cref{sec:ro-freedom,sec:la-coordinates}), and $3N - n_{\mathrm c} - 3$ once the
drift of the centre of mass is removed and kept at zero. A water molecule
with all three of its distances held has $9 - 3 = 6$: the three drifts and
three turnings of \cref{sec:ro-freedom}, and no way left to vibrate. For
64 molecules, 192 atoms, flexible water has $576 - 3 = 573$ freedoms,
water with its 128 O--H bonds held has $573 - 128 = 445$, and rigid water,
with $3\times64 = 192$ distances held, $573 - 192 = 381$.
Chapter~12 needs these counts, since the kinetic energy is shared among
the freedoms.

\begin{keyidea}
RATTLE is velocity Verlet with two corrections along the bonds: after the
drift, the positions are corrected by SHAKE and the half-step velocities
with them; after the second half kick, each relative velocity loses its
part along its bond. Both corrections keep the total momentum; the step
is reversible and symplectic, and agrees with ASE to the rounding of the
arithmetic. Each held distance removes one freedom.
\end{keyidea}

\notebook{11}{correct a pair of velocities, and run RATTLE beside ASE}

\begin{selfcheck}
Atoms of masses 2 and 2, joined along $x$ with $d = 1$, have $\vb{v}_i =
(0.1, 0, 0)$ and $\vb{v}_j = (-0.1, 0.2, 0)$. What are the velocities
after the correction?
\end{selfcheck}
